% Paso 8: Guía de Aplicación Contemporánea
% Proyecto: Análisis Exegético Profundo
% Ministerio "La Gloria es del Señor"

\documentclass[11pt,a4paper]{article}
\usepackage[utf8]{inputenc}
\usepackage[spanish]{babel}
\usepackage[margin=2.5cm]{geometry}
\usepackage{graphicx}
\usepackage{fancyhdr}
\usepackage{xcolor}
\usepackage{tcolorbox}
\usepackage{enumitem}
\usepackage{tabularx}
\usepackage{multirow}
\usepackage{amsfonts}
\usepackage{fontspec}
\usepackage{titlesec}

% Configuración de colores
\definecolor{forestgreen}{RGB}{27,67,50}
\definecolor{sagegreen}{RGB}{45,90,39}
\definecolor{olivegreen}{RGB}{64,145,108}
\definecolor{mintgreen}{RGB}{82,183,136}
\definecolor{lightmint}{RGB}{116,198,157}
\definecolor{softcream}{RGB}{216,243,220}
\definecolor{papergray}{RGB}{118,128,150}

% Corrección para el warning de fancyhdr
\setlength{\headheight}{14pt}

% Configuración de encabezados
\pagestyle{fancy}
\fancyhf{}
\fancyhead[L]{\textcolor{forestgreen}{\textbf{Análisis Exegético Profundo}}}
\fancyhead[R]{\textcolor{papergray}{Paso 8: Aplicación Contemporánea}}
\fancyfoot[C]{\textcolor{papergray}{\thepage}}
\fancyfoot[R]{\textcolor{papergray}{\scriptsize Ministerio "La Gloria es del Señor"}}

% Configuración de títulos
\titleformat{\section}
{\color{forestgreen}\Large\bfseries}
{\thesection}{1em}{}

\titleformat{\subsection}
{\color{sagegreen}\large\bfseries}
{\thesubsection}{1em}{}

% Configuración de cajas
\tcbuselibrary{skins}
\newtcolorbox{infobox}[1]{
    colback=softcream,
    colframe=olivegreen,
    fonttitle=\bfseries,
    title=#1,
    left=10pt,
    right=10pt,
    top=10pt,
    bottom=10pt
}

\newtcolorbox{checklistbox}{
    colback=white,
    colframe=mintgreen,
    left=15pt,
    right=15pt,
    top=10pt,
    bottom=10pt,
    boxrule=2pt
}

\newtcolorbox{applicationbox}{
    colback=lightmint!20,
    colframe=sagegreen,
    left=10pt,
    right=10pt,
    top=8pt,
    bottom=8pt,
    boxrule=1.5pt
}

\newtcolorbox{practicalbox}{
    colback=olivegreen!10,
    colframe=olivegreen,
    left=10pt,
    right=10pt,
    top=8pt,
    bottom=8pt,
    boxrule=1.5pt
}

\newtcolorbox{congregationbox}{
    colback=forestgreen!10,
    colframe=forestgreen,
    left=10pt,
    right=10pt,
    top=8pt,
    bottom=8pt,
    boxrule=1.5pt
}

\begin{document}

% Título principal
\begin{center}
    {\Huge\color{forestgreen}\textbf{PASO 8}}\\[0.5cm]
    {\Large\color{sagegreen}\textbf{APLICACIÓN CONTEMPORÁNEA}}\\[0.3cm]
    {\large\color{papergray}Del Significado Original a la Relevancia Actual}\\[1cm]
    
    \begin{tcolorbox}[colback=olivegreen!10, colframe=olivegreen, width=0.8\textwidth]
        \centering
        \textcolor{forestgreen}{\textbf{Proyecto: Análisis Exegético Profundo}}\\
        \textcolor{papergray}{Metodología de 8 Pasos para Estudio Bíblico Riguroso}
    \end{tcolorbox}
\end{center}

\vspace{1cm}

% Información del estudiante
\begin{infobox}{Información del Proyecto}
\begin{tabularx}{\textwidth}{|X|X|}
\hline
\textbf{Nombre del Estudiante:} & \rule{6cm}{0.4pt} \\[0.8cm]
\hline
\textbf{Texto Bíblico:} & \rule{6cm}{0.4pt} \\[0.8cm]
\hline
\textbf{Congregación/Audiencia:} & \rule{6cm}{0.4pt} \\[0.8cm]
\hline
\textbf{Fecha de Presentación:} & \rule{6cm}{0.4pt} \\[0.8cm]
\hline
\end{tabularx}
\end{infobox}

\section{Objetivo del Paso 8: Aplicación Contemporánea}

La aplicación contemporánea es el puente que conecta el significado original del texto con la vida de la iglesia y el creyente de hoy. Este paso final requiere fidelidad al texto bíblico y relevancia para la audiencia contemporánea, manteniendo el equilibrio entre el "entonces" y el "ahora".

\subsection{Descripción del Proceso}

Desarrollaremos aplicaciones relevantes y legítimas del texto para diferentes audiencias contemporáneas. Identificaremos principios universales, desarrollaremos aplicaciones específicas, y prepararemos el material para la presentación congregacional, asegurándonos de que la aplicación fluya naturalmente de la exégesis.

\begin{infobox}{Principio de Aplicación}
\textbf{La aplicación legítima debe surgir del significado original del texto.} No podemos aplicar lo que el texto no significa, pero debemos aplicar fielmente lo que sí significa en nuestro contexto contemporáneo.
\end{infobox}

\section{Lista de Verificación}

\begin{checklistbox}
\textbf{Marca cada elemento completado:}

\begin{itemize}[leftmargin=1cm]
    \item[$\square$] Identificar principios universales del texto
    \item[$\square$] Desarrollar aplicaciones específicas y concretas
    \item[$\square$] Considerar diferentes audiencias (edades, situaciones)
    \item[$\square$] Verificar legitimidad de las aplicaciones
    \item[$\square$] Preparar ilustraciones contemporáneas relevantes
    \item[$\square$] Desarrollar llamados a la acción específicos
    \item[$\square$] Crear material para diferentes formatos de presentación
    \item[$\square$] Anticipar preguntas y objeciones potenciales
\end{itemize}
\end{checklistbox}

\newpage

\section{Fundamentos para la Aplicación}

\subsection{Resumen del Significado Original}

\begin{applicationbox}
\textbf{Resume brevemente el mensaje principal y los puntos teológicos clave de tu análisis exegético:}
\end{applicationbox}

\vspace{4cm}

\subsection{Principios Universales Identificados}

\textbf{¿Qué principios del texto trascienden la cultura y el tiempo?}

\begin{enumerate}[leftmargin=1cm]
    \item \textbf{Principio Universal \#1:}
    \vspace{2cm}
    
    \item \textbf{Principio Universal \#2:}
    \vspace{2cm}
    
    \item \textbf{Principio Universal \#3:}
    \vspace{2cm}
    
    \item \textbf{Principio Universal \#4:}
    \vspace{2cm}
\end{enumerate}

\subsection{Diferencias Culturales a Considerar}

\textbf{¿Qué aspectos del contexto original son diferentes de nuestro contexto contemporáneo?}

\begin{tabularx}{\textwidth}{|X|X|}
\hline
\textbf{Contexto Original} & \textbf{Contexto Contemporáneo} \\
\hline
& \\[1.5cm]
\hline
& \\[1.5cm]
\hline
& \\[1.5cm]
\hline
& \\[1.5cm]
\hline
\end{tabularx}

\textbf{¿Cómo afectan estas diferencias la forma de aplicar el texto hoy?}
\vspace{3cm}

\newpage

\section{Desarrollo de Aplicaciones Específicas}

\subsection{Aplicación Personal/Individual}

\begin{practicalbox}
\textbf{¿Cómo debe responder el creyente individual a este texto?}
\end{practicalbox}

\textbf{Área: Vida Espiritual Personal}
\vspace{2cm}

\textbf{Área: Relaciones Familiares}
\vspace{2cm}

\textbf{Área: Vida Laboral/Profesional}
\vspace{2cm}

\textbf{Área: Decisiones Éticas}
\vspace{2cm}

\textbf{Área: Carácter Cristiano}
\vspace{2cm}

\subsection{Aplicación Comunitaria/Congregacional}

\begin{congregationbox}
\textbf{¿Cómo debe responder la iglesia como comunidad a este texto?}
\end{congregationbox}

\textbf{Área: Adoración y Vida Litúrgica}
\vspace{2cm}

\textbf{Área: Ministerio y Servicio}
\vspace{2cm}

\textbf{Área: Evangelismo y Misiones}
\vspace{2cm}

\textbf{Área: Vida Comunitaria y Koinonía}
\vspace{2cm}

\textbf{Área: Liderazgo y Gobierno Eclesiástico}
\vspace{2cm}

\newpage

\subsection{Aplicación Social/Cultural}

\textbf{¿Cómo debe la iglesia aplicar este texto en su interacción con la sociedad?}

\textbf{Área: Justicia Social}
\vspace{2cm}

\textbf{Área: Política y Ciudadanía}
\vspace{2cm}

\textbf{Área: Cultura y Arte}
\vspace{2cm}

\textbf{Área: Educación}
\vspace{2cm}

\textbf{Área: Economía y Negocios}
\vspace{2cm}

\section{Aplicaciones por Audiencia Específica}

\subsection{Jóvenes y Adolescentes}

\textbf{¿Cómo es especialmente relevante este texto para los jóvenes de tu congregación?}
\vspace{3cm}

\textbf{¿Qué aplicaciones específicas abordan sus desafíos únicos?}
\vspace{3cm}

\subsection{Adultos/Familias}

\textbf{¿Qué aplicaciones son más relevantes para adultos en diferentes etapas de la vida?}

\textbf{Jóvenes Adultos (20-35 años):}
\vspace{2cm}

\textbf{Adultos de Mediana Edad (35-55 años):}
\vspace{2cm}

\textbf{Adultos Mayores (55+ años):}
\vspace{2cm}

\subsection{Situaciones Especiales}

\textbf{¿Cómo se aplica este texto a personas en circunstancias específicas?}

\textbf{Nuevos Creyentes:}
\vspace{2cm}

\textbf{Creyentes en Crisis:}
\vspace{2cm}

\textbf{Líderes de la Iglesia:}
\vspace{2cm}

\newpage

\section{Verificación de Legitimidad}

\subsection{Evaluación de Aplicaciones}

\textbf{Para cada aplicación principal, verifica:}

\begin{tabular}{|p{4cm}|p{3cm}|p{3cm}|p{3cm}|}
\hline
\textbf{Aplicación} & \textbf{¿Surge del texto?} & \textbf{¿Es contextualmente apropiada?} & \textbf{¿Es prácticamente viable?} \\
\hline
& Sí / No & Sí / No & Sí / No \\[1cm]
\hline
& Sí / No & Sí / No & Sí / No \\[1cm]
\hline
& Sí / No & Sí / No & Sí / No \\[1cm]
\hline
& Sí / No & Sí / No & Sí / No \\[1cm]
\hline
\end{tabular}

\subsection{Errores de Aplicación a Evitar}

\textbf{¿Qué aplicaciones incorrectas del texto has escuchado o podrías anticipar?}
\vspace{3cm}

\textbf{¿Por qué esas aplicaciones serían ilegítimas?}
\vspace{3cm}

\section{Ilustraciones Contemporáneas}

\subsection{Analogías y Metáforas Modernas}

\textbf{¿Qué analogías contemporáneas ayudan a explicar los conceptos del texto?}

\begin{tabularx}{\textwidth}{|X|X|}
\hline
\textbf{Concepto del Texto} & \textbf{Analogía Contemporánea} \\
\hline
& \\[2cm]
\hline
& \\[2cm]
\hline
& \\[2cm]
\hline
\end{tabularx}

\subsection{Ejemplos Contemporáneos}

\textbf{¿Qué ejemplos de la vida real ilustran los principios del texto?}

\textbf{Ejemplo \#1:}
\vspace{2cm}

\textbf{Ejemplo \#2:}
\vspace{2cm}

\textbf{Ejemplo \#3:}
\vspace{2cm}

\newpage

\section{Llamados a la Acción}

\subsection{Respuestas Específicas Requeridas}

\begin{practicalbox}
\textbf{¿Qué acciones concretas debe tomar la audiencia como resultado de este texto?}
\end{practicalbox}

\textbf{Acciones Inmediatas (esta semana):}
\begin{enumerate}[leftmargin=1cm]
    \item \rule{12cm}{0.4pt}
    \item \rule{12cm}{0.4pt}
    \item \rule{12cm}{0.4pt}
\end{enumerate}

\textbf{Compromisos a Mediano Plazo (este mes):}
\begin{enumerate}[leftmargin=1cm]
    \item \rule{12cm}{0.4pt}
    \item \rule{12cm}{0.4pt}
    \item \rule{12cm}{0.4pt}
\end{enumerate}

\textbf{Cambios de Estilo de Vida (a largo plazo):}
\begin{enumerate}[leftmargin=1cm]
    \item \rule{12cm}{0.4pt}
    \item \rule{12cm}{0.4pt}
    \item \rule{12cm}{0.4pt}
\end{enumerate}

\subsection{Preguntas para la Reflexión Personal}

\textbf{¿Qué preguntas ayudarán a la audiencia a aplicar personalmente el texto?}

\begin{enumerate}[leftmargin=1cm]
    \item \rule{12cm}{0.4pt}
    \vspace{0.5cm}
    \item \rule{12cm}{0.4pt}
    \vspace{0.5cm}
    \item \rule{12cm}{0.4pt}
    \vspace{0.5cm}
    \item \rule{12cm}{0.4pt}
    \vspace{0.5cm}
    \item \rule{12cm}{0.4pt}
    \vspace{0.5cm}
\end{enumerate}

\section{Material para la Presentación}

\subsection{Bosquejo de la Presentación}

\textbf{Estructura de tu presentación de 15-20 minutos:}

\textbf{I. Introducción (3-4 minutos)}
\vspace{2cm}

\textbf{II. Análisis del Texto (8-10 minutos)}
\vspace{2cm}

\textbf{III. Mensaje Central (2-3 minutos)}
\vspace{2cm}

\textbf{IV. Aplicación Práctica (3-4 minutos)}
\vspace{2cm}

\textbf{V. Conclusión y Llamado (1-2 minutos)}
\vspace{2cm}

\subsection{Puntos Clave para Recordar}

\textbf{¿Cuáles son los 3 puntos más importantes que tu audiencia debe recordar?}

\begin{enumerate}[leftmargin=1cm]
    \item \rule{12cm}{0.4pt}
    \vspace{1cm}
    \item \rule{12cm}{0.4pt}
    \vspace{1cm}
    \item \rule{12cm}{0.4pt}
    \vspace{1cm}
\end{enumerate}

\newpage

\section{Anticipación de Preguntas}

\subsection{Posibles Preguntas de la Audiencia}

\textbf{¿Qué preguntas podrían surgir durante o después de tu presentación?}

\begin{tabularx}{\textwidth}{|X|X|}
\hline
\textbf{Pregunta Anticipada} & \textbf{Tu Respuesta} \\
\hline
& \\[2cm]
\hline
& \\[2cm]
\hline
& \\[2cm]
\hline
& \\[2cm]
\hline
\end{tabularx}

\subsection{Objeciones Potenciales}

\textbf{¿Qué objeciones o resistencias podrían tener a tus aplicaciones?}
\vspace{3cm}

\textbf{¿Cómo responderías a estas objeciones de manera pastoral y bíblica?}
\vspace{3cm}

\section{Evaluación Final}

\subsection{Revisión de la Aplicación}

\textbf{¿Tus aplicaciones fluyen naturalmente de tu análisis exegético?}
\vspace{2cm}

\textbf{¿Son prácticas y viables para tu audiencia específica?}
\vspace{2cm}

\textbf{¿Mantienen el equilibrio entre gracia y verdad?}
\vspace{2cm}

\textbf{¿Apuntan hacia la transformación por el poder del Espíritu Santo?}
\vspace{2cm}

\subsection{Impacto Esperado}

\textbf{¿Qué esperas que Dios haga en la vida de tu audiencia a través de este texto?}
\vspace{3cm}

\textbf{¿Cómo podrías dar seguimiento para ver el fruto de la aplicación?}
\vspace{3cm}

\begin{infobox}{Recordatorio Final}
La meta de la aplicación no es solo informar sino transformar. Confía en que el Espíritu Santo usará su Palabra fielmente proclamada para obrar en los corazones de tu audiencia.
\end{infobox}

\section{¡Análisis Exegético Completado!}

¡Felicitaciones! Has completado un análisis exegético riguroso y completo. Ahora estás listo para presentar los tesoros de la Palabra de Dios a tu congregación con confianza académica y pasión pastoral.

\vspace{1cm}

\begin{center}
\begin{tcolorbox}[colback=forestgreen!20, colframe=forestgreen, width=0.8\textwidth]
\centering
\textcolor{forestgreen}{\textbf{\Large ¡PROYECTO COMPLETADO!}}\\[0.5cm]
\textcolor{papergray}{\textbf{Has completado exitosamente los 8 pasos del análisis exegético.}}\\
\textcolor{papergray}{Ahora estás equipado para comunicar fielmente la Palabra de Dios.}\\[0.5cm]
\textcolor{sagegreen}{\textit{"Procura con diligencia presentarte a Dios aprobado, como obrero que no tiene de qué avergonzarse, que usa bien la palabra de verdad."}}\\
\textcolor{papergray}{\small 2 Timoteo 2:15}
\end{tcolorbox}
\end{center}

\end{document}